\documentclass[11pt]{article}
\usepackage[T1]{fontenc}
\usepackage[utf8]{inputenc}
\usepackage{amsmath,amssymb,amsthm}
\usepackage[margin=1in]{geometry}
\usepackage{microtype}
\usepackage[hidelinks]{hyperref}
\hypersetup{
  pdftitle={A Roman-Domination Bound for the Graph Burning Number},
  pdfauthor={Xiaojian Zeng}
}

\newtheorem{theorem}{Theorem}
\newtheorem{lemma}[theorem]{Lemma}

\title{A Roman-Domination Bound for the Graph Burning Number}
\author{Xiaojian Zeng\\
PBC School of Finance, Tsinghua University, 100190\\
\texttt{xiaojianzeng@foxmail.com}}
\date{}

\begin{document}
\maketitle

\begin{abstract}
For every finite simple connected graph $G$ with at least two vertices, we
prove
\[
  2b(G)-2\leq \gamma_R(G),
\]
where $b(G)$ is the burning number under the simultaneous-spread convention
and $\gamma_R(G)$ is the Roman domination number.  The proof chooses a
minimum Roman dominating function having as few vertices labelled $1$ as
possible.  The vertices labelled $1$ are then independent and have no
neighbour labelled $2$, which puts them at distance two from the set of
$2$-vertices.  This structure gives a legal burning process of length at most
$\lfloor\gamma_R(G)/2\rfloor+1$.
\end{abstract}

\section{Definitions and context}

All graphs in this note are finite, simple, and undirected.  For a vertex $x$
of a graph $G$, let
\[
  N_r[x]=\{v\in V(G):d_G(v,x)\leq r\}.
\]
Under the simultaneous-spread convention, one previously unburned vertex is
chosen as a new source in each round, and fire already present spreads across
one edge.  Equivalently, the burning number $b(G)$ is the least positive
integer $k$ for which there are vertices $x_1,\ldots,x_k$ such that
\begin{equation}\label{eq:burning}
 V(G)=\bigcup_{i=1}^k N_{k-i}[x_i],
 \qquad d_G(x_i,x_j)\geq j-i\quad(1\leq i<j\leq k).
\end{equation}
This metric formulation and its connection with the spreading process are
standard in the graph-burning literature; see Bessy et
al.~\cite{BessyEtAl2018}.

A \emph{Roman dominating function} (RDF) on $G$ is a map
$f:V(G)\to\{0,1,2\}$ such that every vertex assigned $0$ has a neighbour
assigned $2$.  Its weight is $w(f)=\sum_{v\in V(G)}f(v)$, and the Roman
domination number $\gamma_R(G)$ is the minimum weight of an RDF.  These are
the standard definitions of Cockayne et al.~\cite{CockayneEtAl2004}.

The inequality proved below was recorded as a surviving machine-generated
candidate in the data accompanying AutoGraphForge~\cite{Pastorek2026}.
That system tests candidate graph inequalities on a large finite graph
collection and explicitly distinguishes survival of such tests from a
universal theorem.  We make no claim about priority or about the status of the
inequality before that release; the contribution here is the proof of the
precise universal statement.

\section{The bound}

We first isolate the structure needed from an optimal RDF.

\begin{lemma}\label{lem:rdf}
Let $G$ be connected and have at least two vertices.  There is a minimum RDF
$f$ such that, on writing $V_i=f^{-1}(i)$, the following hold:
\begin{enumerate}
  \item $V_1$ is independent;
  \item there is no edge between $V_1$ and $V_2$;
  \item $V_2$ is nonempty; and
  \item every vertex of $V_1$ is at distance exactly two from $V_2$.
\end{enumerate}
\end{lemma}

\begin{proof}
Among all minimum RDFs, choose $f$ so that $|V_1|$ is as small as possible.
If a vertex of $V_1$ had a neighbour in $V_2$, changing its label from $1$
to $0$ would produce a lighter RDF.  This proves (2).

If adjacent vertices $u,v$ both belonged to $V_1$, change the label of $u$
to $2$ and that of $v$ to $0$.  Every old $0$-vertex keeps its old
$2$-neighbour, while the new $0$-vertex $v$ is defended by $u$.  The weight
is unchanged but there are two fewer $1$-vertices, a contradiction.  Thus
(1) holds.

Suppose that $V_2$ were empty.  Roman feasibility would then force every
vertex to have label $1$.  Choose a vertex $v$ of positive degree, which
exists by connectedness and nontriviality.  Assign $2$ to $v$, assign $0$ to
all its neighbours, and assign $1$ to every other vertex.  The resulting RDF
has weight
\[
  |V(G)|+1-\deg(v)\leq |V(G)|.
\]
Strict inequality contradicts minimum weight, while equality gives a minimum
RDF with fewer $1$-vertices.  Hence (3) holds.

Finally, every vertex of $V_0$ has a neighbour in $V_2$.  A vertex of $V_1$
has a neighbour because $G$ is connected and has at least two vertices.  By
(1) and (2), that neighbour belongs to $V_0$, and consequently it has a
neighbour in $V_2$.  Thus the distance from each $V_1$-vertex to $V_2$ is at
most two and, by (2), is exactly two.  This proves (4).
\end{proof}

\begin{theorem}\label{thm:main}
For every finite simple connected graph $G$ with at least two vertices,
\[
  2b(G)-2\leq \gamma_R(G).
\]
\end{theorem}

\begin{proof}
Choose $f$ as in Lemma~\ref{lem:rdf}, and put
\[
  a=|V_2|,\qquad c=|V_1|.
\]
Then $a\geq1$ and $\gamma_R(G)=2a+c$.

During each of the first $a$ burning rounds, choose an unburned vertex of
$V_2$ whenever one remains; if none remains, choose any unburned vertex.  If
the entire graph has already burned, the desired upper bound is immediate.
Every round in which an unburned $V_2$-vertex remains therefore burns at least
one previously unburned member of $V_2$.  It follows that all of $V_2$ is
burned by the end of round $a$.  Every source is unburned when chosen, so the
construction is legal under \eqref{eq:burning}.

If $c=0$, one additional round lets the fire on $V_2$ reach all of $V_0$,
and hence $b(G)\leq a+1$.  If $c=1$, in round $a+1$ choose the unique
$V_1$-vertex when it is still unburned.  In that same round the existing fire
reaches all of $V_0$, so again $b(G)\leq a+1$.  If $c\geq2$, two additional
rounds let the fire on $V_2$ reach every vertex at distance at most two.
Lemma~\ref{lem:rdf} then gives
\[
  b(G)\leq a+2\leq a+\left\lfloor\frac c2\right\rfloor+1.
\]
The same final bound holds in the first two cases.  Consequently,
\[
  b(G)\leq a+\left\lfloor\frac c2\right\rfloor+1
       =\left\lfloor\frac{\gamma_R(G)}2\right\rfloor+1.
\]
Since the parameters are integral, this implies
$2b(G)-2\leq\gamma_R(G)$.
\end{proof}

The inequality is sharp: for $K_2$, one has $b(K_2)=2$ and
$\gamma_R(K_2)=2$.

\section{Computational check and assistance disclosure}

As a check on the structural transformations and round timing, an exact
enumerator examined all $995$ connected unlabeled graphs of orders two
through seven.  It independently optimized the two parameters, selected the
canonical RDF used in Lemma~\ref{lem:rdf}, and simulated the prescribed
burning strategy; no failure occurred.  The implementation used Python and
NetworkX 3.4.2.  These finite calculations are diagnostic only and are not
used in the proof of Theorem~\ref{thm:main}.  This distinction matters: exact
enumeration can prove a theorem whose domain is itself a specified finite
library, as illustrated by the finite geometric certificates of Zeng et
al.~\cite{ZengEtAl2026}, but it cannot by itself establish the universal
statement considered here.

OpenAI Codex assisted with proof exploration, preparation and execution of
the finite checks, literature triage, adversarial review organization, and
LaTeX editing.  Language-model output and the finite computations are not
treated as mathematical proof; the universal conclusion rests on the
deductive argument above.  The author takes responsibility for the statement
and proof.

\bibliographystyle{plain}
\begin{sloppypar}
\bibliography{references}
\end{sloppypar}

\end{document}
